% Unit 112 terjemahan Bahasa Indonesia dari chapter8.tex baris 1059--1228.
% Sumber: Methods of Algebra, Volume 2, commit
% 9a5803ff77dd3257484cb177f851a73770a59dd3 (CC BY 4.0).
% stable-unit-id: o014.aljabr2.chapter8.homology-computation
% Cakupan: Bagian 8.6 lengkap, Perhitungan Homologi, hingga tepat sebelum
% Bagian 8.7.
% Perbaikan matematis-redaksional terhadap sumber dicatat pada peta segmen:
% koefisien H_0, normalisasi ZS, interval Delta^1, indeks diferensial
% augmentasi, urutan komposisi augmentasi, dan rumus kontraksi E Gamma.

% segment-id: o014.aljabr2.chapter8.homology-computation.h001
\section{Perhitungan Homologi}\label{sec:homology-computation}
\hypertarget{o014.aljabr2.chapter8.homology-computation}{ }
% segment-id: o014.aljabr2.chapter8.homology-computation.p001
Penerapan klasik funktor $\mathrm{C}$ atau $\mathrm{N}$ ialah mendefinisikan
homologi himpunan simplisial, yang berkaitan erat dengan kajian topologi
aljabar.

% segment-id: o014.aljabr2.chapter8.homology-computation.d001
\begin{definition}\label{def:free-simplicial-Ab}
	Untuk setiap $K\in\Obj(\cate{sSet})$, definisikan
	$\Z K\in\Obj(\cate{sAb})$ dengan mengambil $(\Z K)_n$ sebagai
	modul-$\Z$ bebas berbasis $K_n$, sedangkan pemetaan $d_i,s_j$ padanya
	diinduksi oleh pemetaan yang bersesuaian pada $K$. Konstruksi ini
	memberikan funktor $\Z(\cdot):\cate{sSet}\to\cate{sAb}$.
\end{definition}

% segment-id: o014.aljabr2.chapter8.homology-computation.pr001
\begin{proposition}\label{prop:free-simplicial-Ab-adjunction}
	Terdapat adjungsi alami
	% segment-id: o014.aljabr2.chapter8.homology-computation.g001
	$\begin{tikzcd}
		\Z(\cdot):\cate{sSet}\arrow[shift left,r]
		&\cate{sAb}:\text{pelupa}\arrow[shift left,l]
	\end{tikzcd}$.
\end{proposition}
% segment-id: o014.aljabr2.chapter8.homology-computation.pf001
\begin{proof}
	Funktor $\Z(\cdot)$ dan funktor pelupa sama-sama didefinisikan tingkat demi
	tingkat. Seluruh pernyataan langsung mengikuti adjungsi bebas--pelupa
	$\cate{Set}\leftrightarrows\cate{Ab}$.
\end{proof}

% segment-id: o014.aljabr2.chapter8.homology-computation.p002
Karena itu, untuk setiap $X\in\Obj(\cate{sAb})$, memberikan morfisme
$\Z\Delta^n\to X$ sama artinya dengan memberikan suatu unsur $X_n$.

% segment-id: o014.aljabr2.chapter8.homology-computation.p003
Selain itu, $\Z(X\times Y)\simeq\Z X\otimes\Z Y$, dengan $\otimes$ di
ruas kanan dibentuk sebagai $X_n\dotimes{\Z}Y_n$ tingkat demi tingkat.

% segment-id: o014.aljabr2.chapter8.homology-computation.d002
\begin{definition}\label{def:simplicial-set-homology}
	\index[sym1]{Hn}
	\index[sym1]{Hnn}
	Untuk himpunan simplisial $K$ dan setiap $n\in\Z_{\geq0}$, tuliskan
	$\Hm_n(K;\Z):=\Hm_n\bigl(\mathrm{C}(\Z K)\bigr)$; grup ini disebut
	\emph{grup homologi ke-$n$ dari $K$}. Konstruksi ini memberikan keluarga
	funktor $\Hm_n(\cdot;\Z):\cate{sSet}\to\cate{Ab}$.
	% segment-id: o014.aljabr2.chapter8.homology-computation.l001
	\begin{itemize}
		\item Secara lebih umum, untuk modul-$\Z$ $M$, homologi berkoefisien
		$M$ didefinisikan oleh
		$\Hm_n(K;M):=\Hm_n\bigl(\mathrm{C}(\Z K)\otimes M\bigr)$, dengan
		$\mathord\cdot\otimes M$ berarti mengambil
		$\mathord\cdot\dotimes{\Z}M$ pada setiap suku kompleks rantai.

		\item Kohomologi berkoefisien $M$ didefinisikan oleh
		$\Hm^n(K;M):=\Hm^n\bigl(\Hom_{\Z}(\mathrm{C}(\Z K),M)\bigr)$, dengan
		ekspresi di dalam tanda kurung dipandang sebagai kompleks. Berdasarkan
		sifat universal $\Z(\cdot)$, suku ke-$n$ kompleks ini juga dapat
		diidentifikasikan dengan
		% segment-id: o014.aljabr2.chapter8.homology-computation.q001
		\[
		\mathrm{Maps}(K_n,M):=\{\text{pemetaan }K_n\to M\}.
		\]
		Himpunan ini menjadi modul-$\Z$ melalui operasi titik demi titik, dan
		homomorfisme diferensialnya diidentifikasi dengan
		$\sum_i(-1)^i(d_i)^*$.
	\end{itemize}
\end{definition}

% segment-id: o014.aljabr2.chapter8.homology-computation.p004
Jika $\mathrm{C}(\Z K)$ dalam definisi di atas diganti dengan
$\mathrm{N}(\Z K)$, grup homologi (atau kohomologi) yang diperoleh isomorfik
secara kanonik; ini merupakan akibat Teorema~\ref{prop:Dold-Kan-N-C}.
Penggunaan $\mathrm{N}(\Z K)$ kadang lebih sederhana karena
Proposisi~\ref{prop:Dold-Kan-v} memberikan
% segment-id: o014.aljabr2.chapter8.homology-computation.q002
\[
\mathrm{N}(\Z K)_n=\bigoplus_{x\in K_n^{\mathrm{nd}}}\Z x.
\]
Di sini $K_n^{\mathrm{nd}}\subset K_n$ adalah subhimpunan yang terdiri atas
simpleks-$n$ nondegenerat, sedangkan pemetaan
$\mathrm{N}(\Z K)_n\to\mathrm{N}(\Z K)_{n-1}$ diperoleh dengan terlebih
dahulu mengambil $\sum_{i=0}^n(-1)^id_i$, kemudian memproyeksikan ke bagian
$K_{n-1}^{\mathrm{nd}}$. Sebagai contoh,
% segment-id: o014.aljabr2.chapter8.homology-computation.q003
\begin{align*}
	\mathrm{N}(\Z\Delta^0)&=\Z\;\text{yang ditempatkan pada derajat $0$},\\
	\mathrm{C}(\Z\Delta^0)&=
	\left[\cdots\xrightarrow{1}\Z\xrightarrow{0}\Z
	\xrightarrow{1}\Z\xrightarrow{0}\Z\right].
\end{align*}

% segment-id: o014.aljabr2.chapter8.homology-computation.eg001
\begin{example}\label{eg:homology-poset-e}
	Sebagai kasus khusus yang berguna, misalkan diberikan poset kecil takkosong
	$(Q,\leq)$ yang mempunyai unsur terkecil $e$. Ini memberikan pemetaan
	mempertahankan urutan
	% segment-id: o014.aljabr2.chapter8.homology-computation.q004
	\[
	[0]\xrightarrow{0\mapsto e}Q\longrightarrow[0].
	\]
	Pandang $Q$ sebagai kategori dan ambil saraf
	$S:=\mathrm{N}(Q)$ dari Contoh~\ref{eg:nerve-cat}. Berdasarkan definisi
	asli saraf, kita juga mempunyai $\mathrm{N}([n])=\Delta^n$. Dengan
	demikian diperoleh morfisme-morfisme dalam $\cate{sSet}$
	% segment-id: o014.aljabr2.chapter8.homology-computation.q005
	\[
	\Delta^0\xrightarrow{i}S\xrightarrow{q}\Delta^0,
	\qquad qi=\identity_{\Delta^0}.
	\]
	Selanjutnya, diperoleh
	% segment-id: o014.aljabr2.chapter8.homology-computation.q006
	\begin{gather*}
		\Z=\mathrm{N}(\Z\Delta^0)
		\xrightarrow{\mathrm{N}i}\mathrm{N}(\Z S)
		\xrightarrow{\mathrm{N}q}\Z,
	\end{gather*}
	dengan $\Z$ dipandang sebagai kompleks rantai yang ditempatkan pada derajat
	$0$. Akan ditunjukkan bahwa $\mathrm{N}i$ dan $\mathrm{N}q$ saling invers
	dalam kategori homotopi kompleks rantai. Khususnya,
	$\Hm_0(Q;M)\simeq M$, sedangkan
	$\Hm_n(Q;M)=\{0\}$ untuk $n\neq0$.

	Karena $qi=\identity$, cukup dibangun homotopi dari
	$\mathrm{N}i\mathrm{N}q$ ke $\identity$. Untuk setiap $n\geq0$, perhatikan
	bahwa unsur-unsur $S_n^{\mathrm{nd}}$ adalah rantai di $Q$ yang berbentuk
	$q_0<\cdots<q_n$. Definisikan
	% segment-id: o014.aljabr2.chapter8.homology-computation.q007
	\begin{align*}
		h_n:S_n^{\mathrm{nd}}&\longrightarrow
		S_{n+1}^{\mathrm{nd}}\cup\{0\},\qquad n\geq0,\\
		q_0<\cdots<q_n&\longmapsto
		\begin{cases}
			e<q_0<\cdots<q_n,&q_0\neq e,\\
			0,&q_0=e.
		\end{cases}
	\end{align*}
	Perluas pemetaan ini secara linear ke
	$\Z S_n^{\mathrm{nd}}\simeq\mathrm{N}(\Z S)_n$. Pemetaan tersebut langsung
	memenuhi syarat-syarat homotopi; rinciannya diserahkan
	kepada pembaca sebagai latihan.
\end{example}

% segment-id: o014.aljabr2.chapter8.homology-computation.p005
Perhitungan homologi tidak dapat dipisahkan dari homotopi. Di bawah
korespondensi Dold--Kan, homotopi kompleks rantai mencerminkan homotopi objek
simplisial, dan yang terakhir dapat didefinisikan secara kombinatorial.

% segment-id: o014.aljabr2.chapter8.homology-computation.d003
\begin{definition}[Homotopi simplisial]\label{def:simplicial-homotopy}
	\index{homotopi!simplisial (simplicial)}
	Misalkan $X$ dan $Y$ objek simplisial dalam kategori $\mathcal{C}$, serta
	$f,g:X\rightrightarrows Y$ sepasang morfisme. Suatu
	\emph{homotopi simplisial dari $g$ ke $f$} adalah keluarga morfisme
	% segment-id: o014.aljabr2.chapter8.homology-computation.q008
	\[
	h_i=h_i^n:X_n\to Y_{n+1},\qquad0\leq i\leq n,
	\]
	yang secara ringkas ditulis $h$ dan memenuhi
	% segment-id: o014.aljabr2.chapter8.homology-computation.q009
	\begin{align*}
		d_0h_0&=f_n,\\
		d_{n+1}h_n&=g_n,\\
		d_ih_j&=
		\begin{cases}
			h_{j-1}d_i,&0\leq i<j,\\
			d_ih_{i-1},&1\leq i=j,\\
			h_jd_{i-1},&i>j+1,
		\end{cases}\\
		s_ih_j&=
		\begin{cases}
			h_{j+1}s_i,&i\leq j,\\
			h_js_{i-1},&i>j.
		\end{cases}
	\end{align*}
\end{definition}

% Reference: Weibel, Theorem 8.3.12
% segment-id: o014.aljabr2.chapter8.homology-computation.p006
Definisi di atas bersifat kombinatorial. Namun, bila
$\mathcal{C}=\cate{Set}$, definisi itu sama dengan
\eqref{eqn:simplicial-homotopy-prep}, yang didasarkan pada intuisi topologis.
Mengapa demikian? Untuk setiap $n\in\Z_{\geq0}$, tuliskan secara konkret
$(\Delta^1)_n=\{\alpha_{-1},\ldots,\alpha_n\}$, dengan
$\alpha_i:[n]\to[1]$ memetakan bilangan $\leq i$ ke $0$ dan bilangan lainnya
ke $1$. Maka menentukan $H_n:(X\times\Delta^1)_n\to Y_n$ sama artinya dengan
menentukan $H_{-1}^n,\ldots,H_n^n:X_n\to Y_n$.
% segment-id: o014.aljabr2.chapter8.homology-computation.l002
\begin{itemize}
	\item Untuk data homotopi simplisial $(h_i^n)_{i,n}$, definisikan
	$H_{-1}^n=g$, $H_n^n=f$, dan
	$H_i^n=d_{n+1}h_i^n$ untuk $0\leq i\leq n-1$.
	\item Untuk data $H$ dari diagram komutatif
	\eqref{eqn:simplicial-homotopy-prep}, definisikan
	$h_i^n:=H_i^{n+1}s_i:X_n\to Y_{n+1}$.
\end{itemize}
Verifikasi langsung menunjukkan bahwa kedua pemetaan tersebut terdefinisi
dengan baik dan saling invers.

% segment-id: o014.aljabr2.chapter8.homology-computation.p007
Untuk kasus $\mathcal{C}=\cate{Ab}$, gantikan saja
$X\times\Delta^1$ dalam \eqref{eqn:simplicial-homotopy-prep} dengan
$X\otimes\Z\Delta^1$ dan bekerja dalam $\cate{sAb}$; argumen yang sama
berlaku tanpa perubahan. Karena itu, hasil berikut sepenuhnya wajar.

% segment-id: o014.aljabr2.chapter8.homology-computation.pr002
\begin{proposition}
	Tinjau sepasang morfisme $f,g:X\rightrightarrows Y$ dalam
	$\cate{s}\mathcal{A}$. Jika $h$ homotopi simplisial dari $g$ ke $f$, maka
	$\bigl(\sum_{j=0}^n(-1)^jh_j^n\bigr)_{n\geq0}$ mendefinisikan homotopi antara
	$\mathrm{C}f,\mathrm{C}g:\mathrm{C}X\rightrightarrows\mathrm{C}Y$.
\end{proposition}
% segment-id: o014.aljabr2.chapter8.homology-computation.pf002
\begin{proof}
	Seperti biasa, tuliskan semua morfisme muka $X$ dan $Y$ secara seragam
	sebagai $d_i$, dan hilangkan superskrip pada $h_j$; ini tidak menimbulkan
	ambiguitas. Terdapat kesamaan berikut di dalam
	$\Hom_{\mathcal{A}}(X_n,Y_n)$:
	% segment-id: o014.aljabr2.chapter8.homology-computation.q010
	\begin{multline*}
		\sum_{i=0}^{n+1}(-1)^id_i\sum_{j=0}^n(-1)^jh_j
		+\sum_{j=0}^{n-1}(-1)^jh_j\sum_{i=0}^n(-1)^id_i\\
		=f_n-g_n
		+\sum_{\substack{0\leq i\leq n+1\\0\leq j\leq n\\
		(i,j)\neq(0,0),(n+1,n)}}(-1)^{i+j}d_ih_j
		+\sum_{\substack{0\leq i\leq n\\0\leq j\leq n-1}}
		(-1)^{i+j}h_jd_i.
	\end{multline*}
	Jumlah pertama pada baris terakhir dapat ditulis kembali sebagai
	% segment-id: o014.aljabr2.chapter8.homology-computation.q011
	\begin{align*}
		&\sum_{\substack{0\leq i\leq n+1\\i<j\leq n}}
		(-1)^{i+j}h_{j-1}d_i
		+\sum_{i=1}^nd_ih_i-\sum_{i=1}^nd_ih_{i-1}
		+\sum_{\substack{1\leq i\leq n+1\\0\leq j<i-1}}
		(-1)^{i+j}h_jd_{i-1}\\
		&\qquad=-\sum_{\substack{0\leq i\leq n\\i\leq j\leq n-1}}
		(-1)^{i+j}h_jd_i
		-\sum_{\substack{0\leq i\leq n\\0\leq j<i}}
		(-1)^{i+j}h_jd_i.
	\end{align*}
	Di sini digunakan
	$1\leq i\leq n\implies d_ih_i=d_ih_{i-1}$. Ekspresi ini membatalkan
	jumlah kedua dan hanya menyisakan $f_n-g_n$.
\end{proof}

% segment-id: o014.aljabr2.chapter8.homology-computation.p008
Selanjutnya kita meninjau objek semisimplisial teraugmentasi
(Catatan~\ref{rem:semisimplicial}) dan kompleks rantai yang bersesuaian;
hasil-hasil ini akan digunakan dalam \S\ref{sec:bar-resolution}. Kecuali
dinyatakan lain, morfisme augmentasi objek semisimplisial teraugmentasi $X$
selalu ditulis sebagai $\epsilon:X_0\to X_{-1}$.

% segment-id: o014.aljabr2.chapter8.homology-computation.d004
\begin{definition}\label{def:augmented-C-cplx}
	\index[sym1]{CXaug@$\mathrm{C}^{\mathrm{aug}} X$}
	Untuk objek semisimplisial teraugmentasi $X$ dalam kategori aditif
	$\mathcal{A}$, kompleks $\mathrm{C}X$ yang bersesuaian dapat diperluas
	menjadi kompleks rantai
	% segment-id: o014.aljabr2.chapter8.homology-computation.q012
	\begin{equation*}
		\begin{gathered}
			\mathrm{C}^{\mathrm{aug}}X:=
			\left[\cdots\xrightarrow{\partial_3}X_2
			\xrightarrow{\partial_2}X_1
			\xrightarrow{\partial_1}X_0
			\xrightarrow{\partial_0:=\epsilon}X_{-1}
			\to0\to\cdots\right],\\
			\partial_n:=\sum_{i=0}^n(-1)^id_i:X_n\to X_{n-1},
			\qquad n\geq1.
		\end{gathered}
	\end{equation*}
\end{definition}

% segment-id: o014.aljabr2.chapter8.homology-computation.p009
Dari $\epsilon d_0=\epsilon d_1$ terlihat bahwa
$\mathrm{C}^{\mathrm{aug}}X$ memang objek
$\cate{Ch}_{\geq-1}(\mathcal{A})$. Kompleks ini juga dapat dipandang sebagai
morfisme dalam $\cate{Ch}_{\geq0}(\mathcal{A})$
% segment-id: o014.aljabr2.chapter8.homology-computation.q013
\[
\epsilon:\mathrm{C}X\to
\underbracket{X_{-1}}_{\text{ditempatkan pada derajat $0$}},
\qquad\text{$\epsilon$ pada derajat $0$, dan $0$ pada derajat lainnya}.
\]

% segment-id: o014.aljabr2.chapter8.homology-computation.d005
\begin{definition}\label{def:contractible-aug}
	\index{dapat dikontraksi (contractible)}
	Misalkan $X$ objek semisimplisial teraugmentasi dalam sebarang kategori
	$\mathcal{C}$. Jika terdapat keluarga morfisme
	$k_n:X_n\to X_{n+1}$, dengan $n\in\Z_{\geq-1}$, gunakan istilah
	\emph{dapat dikontraksi dari kiri} atau \emph{dapat dikontraksi dari
	kanan} bagi $X$, sesuai dengan kolom syarat yang dipenuhi di bawah ini:
	% segment-id: o014.aljabr2.chapter8.homology-computation.tb001
	\begin{center}
	\begin{tabular}{|c|c|}\hline
		Dapat dikontraksi dari kiri&Dapat dikontraksi dari kanan\\\hline
		$\epsilon k_{-1}=\identity_{X_{-1}}$
		&$\epsilon k_{-1}=\identity_{X_{-1}}$\\
		$d_0k_n=\identity_{X_n}$
		&$d_{n+1}k_n=\identity_{X_n}$\\
		$d_ik_n=k_{n-1}d_{i-1}\;(1\leq i\leq n+1,\;n\geq1)$
		&$d_ik_n=k_{n-1}d_i\;(0\leq i\leq n,\;n\geq1)$\\
		$d_1k_0=k_{-1}\epsilon$
		&$d_0k_0=k_{-1}\epsilon$\\\hline
	\end{tabular}
	\end{center}
\end{definition}

% segment-id: o014.aljabr2.chapter8.homology-computation.p010
Periksa bahwa sifat dapat dikontraksi dari kiri dan kanan
saling dual melalui pembalikan urutan dalam pengertian
Catatan~\ref{rem:simpDelta-order-reversal}. Selain itu, keduanya merupakan
sifat ``absolut'': jika $X$ dapat dikontraksi dari kiri (atau kanan), dengan
data yang bersesuaian $(k_n)_{n\geq-1}$, maka untuk setiap funktor
$F:\mathcal{C}\to\mathcal{D}$, data $(Fk_n)_{n\geq-1}$ membuat $F(X)$
dapat dikontraksi dari kiri (atau kanan).

% segment-id: o014.aljabr2.chapter8.homology-computation.eg002
\begin{example}\label{eg:classifying-contractible}
	Untuk sembarang monoid $\Gamma$, jadikan $\mathrm{E}\Gamma$ dari
	Contoh~\ref{eg:classifying-space} himpunan simplisial teraugmentasi dengan
	$(\mathrm{E}\Gamma)_{-1}=\{1\}$ dan morfisme augmentasi
	$\epsilon:(\mathrm{E}\Gamma)_0\to(\mathrm{E}\Gamma)_{-1}$ berupa pemetaan
	konstan. Objek ini dapat dikontraksi dari kanan: untuk $n\geq-1$,
	definisikan
	$k_n:(\mathrm{E}\Gamma)_n\to(\mathrm{E}\Gamma)_{n+1}$ melalui
	$(g_1,\ldots,g_{n+1})\mapsto(g_1,\ldots,g_{n+1},1)$.
	Pemetaan ini jelas memenuhi $\epsilon k_{-1}=\identity$; setelah
	definisi diuraikan, pemeriksaan bahwa
	$d_{n+1}k_n=\identity$,
	$d_ik_n=k_{n-1}d_i$ untuk $0\leq i\leq n$, $n\geq1$, dan
	$d_0k_0=k_{-1}\epsilon$ juga langsung.

	Jika $\Bbbk$ gelanggang komutatif, identifikasikan kompleks rantai
	$\mathrm{C}(\Bbbk\mathrm{E}\Gamma)$ dengan kompleks rantai
	$\mathsf{L}=(\mathsf{L}_n,\partial'_n)_n$ dari
	Definisi~\ref{def:resolution-trivial-module}, sebagaimana dalam
	Contoh~\ref{eg:classifying-space-std-cplx}. Augmentasi
	$\mathsf{L}\to\Bbbk$ persis merupakan augmentasi resolusi dalam
	Proposisi~\ref{prop:trivial-mod-resolution}. Menurut proposisi berikut,
	kontraktibilitas $\mathrm{E}\Gamma$ (dan dengan
	demikian $\Bbbk\mathrm{E}\Gamma$) memberikan penjelasan topologis bagi
	fakta bahwa $\mathsf{L}\to\Bbbk$ adalah kuasi-isomorfisme.
\end{example}

% segment-id: o014.aljabr2.chapter8.homology-computation.pr003
\begin{proposition}\label{prop:contractible-null-homotopic}
	Misalkan $\mathcal{A}$ kategori aditif dan $X$ objek semisimplisial
	teraugmentasi dalam $\mathcal{A}$ yang dapat dikontraksi dari kiri atau
	kanan. Maka morfisme identitas dari
	$\mathrm{C}^{\mathrm{aug}}X$ ke dirinya sendiri homotopik terhadap nol.
\end{proposition}
% segment-id: o014.aljabr2.chapter8.homology-computation.pf003
\begin{proof}
	Pertama, tinjau kasus dapat dikontraksi dari kiri dan ambil keluarga
	morfisme $k_n:X_n\to X_{n+1}$ dalam definisinya. Definisikan $k_n=0$ untuk
	$n<-1$ (atau $\partial_n=0$ untuk $n<0$), dan definisikan pula
	$\partial_0=\epsilon$. Kita hendak menggunakan $(k_n)_n$ sebagai saksi
	homotopi nol, yakni membuktikan
	% segment-id: o014.aljabr2.chapter8.homology-computation.q014
	\[
	\partial_{n+1}k_n+k_{n-1}\partial_n=\identity_{X_n}.
	\]

	Untuk $n<-1$, kedua ruas sama dengan $0$. Untuk $n=-1$, ruas kiri sama
	dengan $\identity_{X_{-1}}$. Untuk $n=0$, ruas kiri ialah
	$d_0k_0-d_1k_0+k_{-1}\epsilon=\identity_{X_0}$. Untuk $n\geq1$,
	perhitungan ruas kiri memberikan
	% segment-id: o014.aljabr2.chapter8.homology-computation.q015
	\begin{multline*}
		\sum_{i=0}^{n+1}(-1)^id_ik_n
		+\sum_{i=0}^n(-1)^ik_{n-1}d_i\\
		=\identity_{X_n}
		+\sum_{i=1}^{n+1}(-1)^ik_{n-1}d_{i-1}
		+\sum_{i=0}^n(-1)^ik_{n-1}d_i
		=\identity_{X_n}.
	\end{multline*}

	Untuk kasus dapat dikontraksi dari kanan, cukup gunakan
	$\bigl((-1)^{n+1}k_n\bigr)_n$ sebagai gantinya.
\end{proof}
